\section{The complete rank-four descent and its downstream use}
\label{app:rank-proof}

This appendix specifies the local completeness and comparison statements
used in Theorem~\ref{thm:rational-t}.  Write $K=K_7=\Q(\sqrt{-7})$ and
$J=J_+$.  The points and curve are those of Appendix~\ref{app:Cplus}.
Put
\[
 H_0=\langle D_2,D_3,D_4,D_5\rangle,\qquad
 H_1=\langle D_2,E,D_4,D_5\rangle.
\]
The relation \eqref{eq:Erelation} implies $H_0\subseteq H_1$.
We distinguish these subgroups throughout: the $2$-saturation concerns
$H_0$, while the $5$-saturation concerns $H_1$.

\subsection{The containing space and complete local conditions}

The global class-group, unit and support calculations give an ordered
$66$-dimensional space of squareclass representatives containing every
relevant fake Selmer class.  The bitangent and syzygetic-tetrad construction
identifies the actual descent modules with those of
\cite[Corollary~12.5 and Appendix~A]{BPS}.  This identification uses the
exact line/contact identities, all $315$ syzygetic tetrads, the rational
even theta characteristic, and the seven blocks supplied by the exact
field tower.  The resulting Galois action is contained in
$\operatorname{AGL}_3(\F_2)$ and is transitive on the $28$ bitangents;
no equality with the full affine group is needed.

For each selected place, the restriction matrix is reconstructed by
arithmetic in the actual local bitangent algebra.  The allowed subspace
contains the entire local Kummer image.  At four split good primes and the
nonsplit prime $1051$, this is proved by exhibiting sufficiently many
independent point images and using the local two-torsion bound.  The same
argument gives image dimensions two and three at the places above $3$ and
$5$.  At $7$ and $2$, the fake-kernel corrections described below supply
the required stopping bounds.  The successive intersections have dimensions
\begin{equation}\label{eq:rank-dimensions}
 66\longrightarrow23\longrightarrow22\longrightarrow20
 \longrightarrow19\longrightarrow18\longrightarrow15\longrightarrow12.
\end{equation}
The stages are, respectively, the four split primes, the nonsplit prime,
the places above $3,5,7$, and the two dyadic places.  The other archived
prime conditions and the norm condition are redundant on this route.
Each of the $66$ restriction columns at the selected places is checked.
Independent linear algebra verifies the intersections and their equality
with the arrays consumed by the subsequent stages.

The diagonal subspace has dimension seven.  If $Q$ denotes the resulting
quotient and $B$ the fake image of $H_0$, exact coordinate identities give
\begin{equation}\label{eq:rank-containing-quotient}
 Q=B\oplus\langle c\rangle,\qquad \dim B=4,\qquad \dim Q=5.
\end{equation}
Here $c$ is the same literal squareclass used in the $5$-adic obstruction
below.  Equality in \eqref{eq:rank-containing-quotient} is checked directly;
it is not inferred from dimensions alone.  It is enough that $Q$ contains
the relevant fake Selmer image: no assertion that every survivor is a
Selmer class is used.

\subsection{The local stopping argument at seven}

The actual decomposition data imply
$\dim_{\F_2}J[2](K_v)=2$ at $v\mid7$, hence
$\#J(K_v)/2J(K_v)=4$.  Set $z_3=[P_4]-[P_1]$.  Exact local square tests
at all seven components of the bitangent algebra give
\[
                      f(z_3)=3\lambda^2.
\]
Thus $z_3$ belongs to the full fake kernel.

A fresh computation in the degree-$42$ target algebra uses a completion
$A/K_v$ with $(e,f)=(4,1)$.  Its incidence algebra is the product of two
unramified quadratic extensions of $A$.  An exact integral Eisenstein
model over $\Q_7$ has degree eight.  Strong Hensel bounds, followed by
Krasner's lemma, bind its approximate generator to the exact target root.
Gauss valuation bounds propagate that root error through every rational
line, incidence and conic coefficient.  The incidence factors are lifted
from coprime factors; a single uniformizer change in one factor reveals
the second irreducible residual quadratic.

After removing the even valuations of $f(z_3)/3$, the two unit square roots
exist by the simple-root Hensel criterion over $\F_{49}$.  Their norm
residues are both one, while the normalized conic quotient has residue
six.  Since every target tetrad has weight four, $\theta(3)=3^2$, and
\cite[Lemma~A.24(a)]{BPS} gives
\[
 \zeta(z_3)=\frac{r(z_3)}{\tau_*(\lambda)\,3^2}=-1.
\]
The exact contact identity makes this quotient an element of $\mu_2$;
reduction in odd residue characteristic distinguishes its two possible
values.  At working precisions $100$ and $120$ over $\Q_7$, strengthened
factor lifting gives separation margins $255$ and $511$ in $A$.
The distinct cycle $[P_2]-[P_1]$, also with diagonal $3$, gives $+1$.
Changing either square-root sign leaves its quadratic norm unchanged.

Every local affine-group candidate with the measured bitangent and
incidence profiles is checked: evaluation on this target descends to
$R^\vee(K_v)/qE^\vee(K_v)$.  Some candidate characters vanish, so the
profile alone is not a nonvanishing proof.  The actual value $-1$ proves
that the physical correction is nonzero.

A genuine closed point over the unramified quadratic extension supplies
a nonzero fake image.  Its local model is transported to the printed
quartic by
\[
                         Z_{\rm original}=-7Z_{\rm local}.
\]
The point satisfies a strong Hensel inequality, all seven line-value
squareclasses are stable, and its norm-line vector equals the archived
allowed line modulo the two diagonal classes.  There is therefore both a
nonzero fake image and a nonzero fake kernel in a group of order four.
The stopping equality of \cite[Lemma~A.24(b)]{BPS} proves that the full
fake image has dimension one, as required in
\eqref{eq:rank-dimensions}.

\subsection{The second dyadic correction}

At the dyadic place where $w=(1+s)/2\equiv0\pmod2$, the local Kummer
group has dimension $3+\dim J[2](K_v)=4$.  Exhaustive local square tests
prove that the four known point differences have fake-image rank two.
The resolved point $C_0$ gives an independent third image.  Two further
actual local cycles have fake-kernel relations
\[
 [\text{old mask},C_0\text{ coefficient},\text{diagonal}]
             =[1,1,-10],\qquad [1,0,2].
\]
The exact degree-six target fields and their quadratic incidence factors
are certified by their local integral bases, ramification and residue
polynomials.  In each factor a square-root witness is certified by Hensel
lifting.  No uniqueness claim about the finite root search is needed.
Both cycles give correction profile $(0,1)$ on the two sextic carriers;
a stronger working accuracy gives the same profile.

For all $28$ affine-group candidates compatible with the actual target
splitting $[3,3,6,6,24]$, the two sextic evaluations annihilate
$qE^\vee(K_v)$ and form independent characters of the two-dimensional
correction quotient.  The nonzero entry therefore detects a genuine
nonzero fake-kernel correction.  The three-dimensional fake image and
this nonzero kernel fill the local Kummer group of dimension four, so the
full fake image has dimension three.  Together with the independently
verified first-dyadic image, this proves the last two cuts in
\eqref{eq:rank-dimensions}.

\subsection{The complementary class and the rank conclusion}

The global correction module is one-dimensional.  In the actual
$5$-adic correction quotient, which has dimension two, the two coherent
global lifts of the literal $c$ in \eqref{eq:rank-containing-quotient}
have obstruction vectors
\[
 \omega_{00}=(1,0),\qquad \omega_{11}=(0,1),\qquad
                         \kappa_5(\rho)=(1,1).
\]
Here $\rho$ generates the global correction module; the last vector is
the difference of the two lift obstructions, not either obstruction
itself.  The functional $(u,v)\mapsto u+v$ kills this difference and is
nonzero on both lifts.  The true-to-fake Selmer comparison kernel is zero,
because this global correction generator restricts nontrivially and the
actual local fake map at $5$ has trivial kernel.

If a true Selmer class had fake image $c+b$ with $b\in B$, subtracting
the Kummer class of a point of $H_0$ mapping to $b$ would give an
everywhere locally Kummer lift of $c$, contrary to the displayed
obstructions.  Thus the true Selmer image lies in $B$, and its dimension
is at most four.  The independently verified lower bound is four.  Since
$J(K)[2]=0$, the Kummer exact sequence gives
\[
 \dim_{\F_2}\operatorname{Sel}^{2}(J/K)=\rank J(K)=4,
                      \qquad \Sha(J/K)[2]=0.
\]
The direct comparison of literal squareclasses, exact target norm roots,
coherent sign relations, and actual local images are included in the
release; cached obstruction bits alone are not used as proof.

Finally, the independent $2$-saturation gives $[J(K):H_0]$ odd.  The
relation \eqref{eq:Erelation} and the checked $5$-saturation of $H_1$
give
\[
                         \gcd([J(K):H_1],400)=1.
\]
Consequently $H_1\to J(K)/400J(K)$ is surjective.  This is the coefficient
coverage used by the finite sieve; full saturation at every prime is not
required.  A characteristic-zero logarithmic functional vanishing on
$H_0$ vanishes on $J(K)$, since $H_0$ has finite index.  Thus the verified
logarithmic annihilator used in Theorem~\ref{thm:rational-t} is global.
